\newcommand{\nomeEstado}{BAHIA}
\newcommand{\siglaEstado}{BA}
\newcommand{\nomeConcessionaria}{NEOENERGIA BAHIA}
\newcommand{\cnpjConcessionaria}{15.139.629/0001-94}
\newcommand{\termoCIP}{Sexto}
\newcommand{\presidenteConcessionaria}{Sr. Thiago Guth.}

\newcommand{\empresaResponsavel}{ABEL}
\newcommand{\nomeMunicipio}{Maiquinique}
\newcommand{\nomeMunicipioMaisculo}{\MakeUppercase{\nomeMunicipio}}
\newcommand{\IPestimada}{3179701}
\newcommand{\cnpjMunicipio}{13.751.821/0001-01}
\newcommand{\sedePrefeitura}{Rua Francisco Martins, nº 01, Centro, Maiquinique – BA, CEP 45770-000}
\newcommand{\nomePrefeito}{Valéria Ferreira Silveira Moreira}
\newcommand{\generoPrefeitura}{pela Prefeita}
\newcommand{\numeroContrato}{0092/2024}
\newcommand{\quantidadeAditivos}{1}
\newcommand{\legitimidade}{1}
\newcommand{\publicacao}{Diário Oficial do Município}
\newcommand{\dataPublicacao}{20 de Junho de 2024}
\newcommand{\tituloClausula}{Cláusula Segunda – Do Objeto}
\newcommand{\clausula}{Deverá a empresa efetuar levantamento dos créditos a que faz jus o município, referentes aos pagamentos indevidos à concessionária de energia elétrica, conforme os prazos estabelecidos na Resolução Normativa da ANEEL nº 1.000, de 7 de dezembro de 2021, especialmente o art. 323, §2º, inciso I, e suas alterações posteriores, podendo para tanto, representar administrativamente o Município de Maiquinique - BA perante a Agência Nacional de Energia Elétrica – ANEEL, em todos os processos, reclamações, recursos e demais atos decorrentes de suas competências.}
\newcommand{\telefone}{(77) 3275-2179}
\newcommand{\email}{prefeitura@maiquinique.ba.gov.br}
